\section{Conclusion}
\label{sec:conclusion}

We have shown that better per-passage retrieval can make answer
faithfulness \emph{worse} when refinement fragments the coherence of
the retrieved set. The phenomenon is robust across $5$ datasets, $5$
generator families ($7$B--$120$B), $9$ ablation axes, and $7{,}000+$
queries. We have introduced a generator-free retrieval-time signal
(CCS) that detects the failure, validated it cross-signal against the
model's own self-confidence, and shipped it as a deployment policy
(HCPC-v$2$) that recovers faithfulness with $0\%$ hallucination at one
pairwise-similarity computation per query. The benchmark, the
$\texttt{pip install}$-able metric, the LangChain integration, and a
permanent Zenodo DOI are released.

The broader message is methodological. RAG evaluation has been
optimising the wrong quantity. We have been measuring how good each
passage looks on its own, when what the generator actually needs is
evidence that fits together. CCS is a small step toward measuring the
right thing.

\paragraph{Broader impact.}
RAG is increasingly deployed in high-stakes settings (medical QA,
legal research, customer support) where unsupported claims have
direct downstream cost. Our work identifies a counter-intuitive
failure mode in the dominant deployment pattern (rerank +
per-passage refinement) and provides a generator-free,
deployment-time signal that flags it. The expected positive impact
is on deployers who can now flag low-coherence retrieval sets before
they reach the generator. A possible negative impact is mis-use of
the diagnostic to over-confidently certify ``high-coherence'' outputs
as faithful: we explicitly mark CCS as a \emph{predictive signal,
not a guarantee}, in \S\ref{sec:discussion:where_ccs_fails} and
identify three regimes where it does not predict.
